\documentclass[../../../include/open-logic-section]{subfiles}

\begin{document}

\olfileid{sth}{spine}{zf}
\olsection{$\Z$ و~$\ZF$: نقطهٔ عطفی دیگر}

با بنیاد به نقطهٔ عطف مهم دیگری می‌رسیم. نظریه‌های~$\Zminus$ و~$\ZFminus$
را بررسی کردیم و گفتیم آن‌ها نظریه‌هایی هستند که چیزی معین را «کم» دارند.
آن چیز بنیاد است. پس:

\begin{defn}
نظریهٔ~$\Z$ بنیاد را به~$\Zminus$ می‌افزاید. بنابراین اصول موضوع آن
عبارت‌اند از گسترش‌مندی، اجتماع، زوج‌ها، مجموعه‌های توانی، نامتناهی،
بنیاد و همهٔ نمونه‌های طرحوارهٔ جداسازی.

نظریهٔ~$\ZF$ بنیاد را به~$\ZFminus$ می‌افزاید. به‌بیانی دیگر، $\ZF$
همهٔ نمونه‌های جایگزینی را به~$\Z$ می‌افزاید.
\end{defn}

بااین‌حال، شاید پرسشی برایتان پیش آمده باشد. اگر قاعده‌مندی نسبت
به~$\ZFminus$ با بنیاد هم‌ارز است و توجیه قاعده‌مندی روشن است، چرا این
راه طولانی را طی کنیم و به‌جای قاعده‌مندی، بنیاد را اصل موضوع پایهٔ خود
قرار دهیم؟

اگر دلایل تاریخی را کنار بگذاریم (اینکه چه کسی چه چیزی را چه هنگام
صورت‌بندی کرد)، دلیل اصلی این است که بنیاد را می‌توان بدون استفاده از
تعریف~$V_\alpha$ها بیان کرد. آن تعریف بر همهٔ کارهای
\olref[recursion]{sec} متکی بود: برای نشان‌دادن موجه‌بودنش باید بازگشت
ترامتناهی را اثبات می‌کردیم. اما اثبات بازگشت ترامتناهی از
\emph{جایگزینی} استفاده کرد. بنابراین، هرچند بنیاد و قاعده‌مندی نسبت
به~$\ZFminus$ هم‌ارزند، نسبت به~$\Zminus$ هم‌ارز نیستند.

در واقع، وضع از آنچه این ملاحظهٔ ساده نشان می‌دهد جدی‌تر است. هرچند این
موضوع بسیار فراتر از دامنهٔ این کتاب است، معلوم می‌شود هم~$\Zminus$ و
هم~$\Z$ برای تعریف~$V_\alpha$ها بیش از حد ضعیف‌اند. پس اگر تنها
در~$\Z$ کار می‌کنید، قاعده‌مندی (با صورت‌بندی ما) حتی \emph{معنا} هم
ندارد. از همین رو اصل موضوع رسمی ما بنیاد است، نه قاعده‌مندی.

از این پس، مگر آنکه خلافش گفته شود، بدون توضیح بیشتر در~$\ZF$ کار خواهیم
کرد.

\end{document}
